\begin{summarybox}
\textbf{一句话核心：先看方程是否降次，再用韦达，最后判断实点。}

\textbf{横开口速查。} 对 $x^2/a^2-y^2/b^2=1$，记
$M=B^2b^2-A^2a^2$、$H=C^2+M$。
\[
 \begin{aligned}
 x_1+x_2&=\frac{2ACa^2}{M},&
 x_1x_2&=-\frac{a^2(C^2+B^2b^2)}{M},\\
 y_1+y_2&=-\frac{2BCb^2}{M},&
 y_1y_2&=\frac{b^2(C^2-A^2a^2)}{M}.
 \end{aligned}
\]
消去 $y$ 需 $B\ne0$；消去 $x$ 需 $A\ne0$。二次韦达须 $M\ne0$。
\[
 \begin{array}{c|c}
 \text{条件}&\text{几何结果}\\\hline
 M\ne0,\ H>0&\text{两个不同交点}\\
 M\ne0,\ H=0&\text{相切}\\
 M\ne0,\ H<0&\text{无实交点}\\
 M=0,\ C\ne0&\text{一个交点，非相切}\\
 M=0,\ C=0&\text{渐近线，无有限交点}
 \end{array}
\]
\textbf{真正有弦时：}
\[
 S=\left(\frac{ACa^2}{M},-\frac{BCb^2}{M}\right),\qquad
 PQ=\frac{2ab\sqrt{(A^2+B^2)H}}{|M|}.
\]
\textbf{纵开口不用重背：} 同时交换 $x,y$ 和 $A,B$，把 $M$ 换成 $N=A^2b^2-B^2a^2$，把 $H$ 换成 $K=C^2+N$。

\textbf{最后四查。}
\begin{enumerate}[leftmargin=*,itemsep=.4em]
 \item 标准化：$a^2,b^2$ 是分母；正号所在的轴决定开口方向。
 \item 消元合法：不除以零；二次项为零先回原方程。
 \item 几何存在：不要把一次交点当切点，也不要把复根当实交点。
 \item 长度完整：根差只是投影；分母取绝对值，竖直线换纵坐标。
\end{enumerate}
\end{summarybox}
% ===== 新增：考场两点法速记（原总结完整保留）=====
\begin{summarybox}
\textbf{新增：考场最低门槛速记——减法求和，交叉乘减求积。}
对横开口 $\dfrac{x^2}{a^2}-\dfrac{y^2}{b^2}=1$ 和 $y=kx+m$，
设有两个不同实交点，且 $b^2-a^2k^2\ne0$。

\textbf{一、求和（两个曲线方程直接相减）：}
\[
 \frac{x_1+x_2}{a^2}=\frac{k(y_1+y_2)}{b^2},\qquad
 y_1+y_2=k(x_1+x_2)+2m
\]
因此
\[
 \boxed{\left(\frac1{a^2}-\frac{k^2}{b^2}\right)(x_1+x_2)
 =\frac{2km}{b^2}}.
\]
\textbf{二、求积（第一个方程乘 $x_2$，第二个乘 $x_1$，再相减）：}
一次交叉项抵消，因此
\[
 \boxed{\left(\frac1{a^2}-\frac{k^2}{b^2}\right)x_1x_2
 =-1-\frac{m^2}{b^2}}.
\]
\textbf{口令：}\fbox{直接减，出和；交叉乘减，出积。}
这是原来韦达公式的\textbf{低记忆负担推导法}：不用背 $M$ 的长分式，
但必须先确认有两个不同实交点；竖直线及二次项退化的情形仍按原方法处理。
\end{summarybox}

% ===== 再补充：两点法写题时的四行自检卡 =====
\begin{summarybox}
\textbf{新增：把推导过程压缩为“四行考场自检”。}
\[
 \frac{x^2}{a^2}-\frac{y^2}{b^2}=1,\qquad y=kx+m.
\]
\textbf{求和：} 两点双曲线式相减 $\to$ 平方差
$\to$ 用 $y_1-y_2=k(x_1-x_2)$ 约分
$\to$ 用 $y_1+y_2=k(x_1+x_2)+2m$：
\[
 \boxed{\left(\frac1{a^2}-\frac{k^2}{b^2}\right)(x_1+x_2)
 =\frac{2km}{b^2}}.
\]
\textbf{求积：} 设 $F(x)=x^2/a^2-(kx+m)^2/b^2-1$，
用 $x_2F(x_1)-x_1F(x_2)=0$：
\[
 \boxed{\left(\frac1{a^2}-\frac{k^2}{b^2}\right)x_1x_2
 =-1-\frac{m^2}{b^2}}.
\]
\textbf{最关键的区别：} 求和利用\emph{平方差}；
求积利用\emph{一次项抵消}。
以上作为除法推导要求 $x_1\ne x_2$，且方程二次项非零；
竖直直线与退化情形仍按前面的条件检查。
\end{summarybox}
